% =============================================================================
% DocumentoStudio1 — Guida allo studio del software asa-deliveroo-agent
% Documento interno (non è il report d'esame): serve a comprendere il
% software in ogni sua parte per poterlo evolvere.
% Compilare con: pdflatex DocumentoStudio1.tex  (x2 per indice e riferimenti)
% =============================================================================
\documentclass[11pt,a4paper]{article}

\usepackage[utf8]{inputenc}
\usepackage[T1]{fontenc}
\usepackage[italian]{babel}
\usepackage[margin=2.4cm]{geometry}
\usepackage{amsmath}
\usepackage{booktabs}
\usepackage{longtable}
\usepackage{xcolor}
\usepackage{listings}
\usepackage[colorlinks=true, linkcolor=blue!60!black, urlcolor=blue!60!black]{hyperref}

% --- stile per i frammenti JavaScript -------------------------------------
\lstdefinelanguage{JavaScript}{
  keywords={break,case,catch,class,const,continue,default,delete,do,else,
    export,extends,false,finally,for,function,if,import,in,instanceof,let,
    new,null,return,static,super,switch,this,throw,true,try,typeof,var,
    void,while,with,async,await,get,set,of},
  sensitive=true,
  comment=[l]{//},
  morecomment=[s]{/*}{*/},
  morestring=[b]',
  morestring=[b]",
  morestring=[b]`
}
\lstset{
  language=JavaScript,
  basicstyle=\ttfamily\small,
  keywordstyle=\color{blue!60!black}\bfseries,
  commentstyle=\color{green!40!black}\itshape,
  stringstyle=\color{red!50!black},
  breaklines=true,
  frame=single,
  framerule=0.3pt,
  columns=fullflexible,
  keepspaces=true,
  showstringspaces=false,
}

\newcommand{\file}[1]{\texttt{#1}}
\newcommand{\code}[1]{\texttt{#1}}

\title{\textbf{DocumentoStudio1}\\[2mm]
  \large Guida allo studio di \texttt{asa-deliveroo-agent}\\
  Architettura, stato attuale e direzioni di evoluzione}
\author{Progetto d'esame ASA --- UniTn\\
  \small Repository: \url{https://github.com/Giambac/asa-deliveroo-agent}}
\date{11 giugno 2026}

\begin{document}

\maketitle

\tableofcontents
\newpage

% ===========================================================================
\section{Scopo di questo documento}
% ===========================================================================

Questo \emph{non} è il report d'esame (che vive in \file{report/main.tex} e
va scritto in inglese a progetto finito). È una guida di studio in italiano
del software costruito finora: spiega \textbf{che cosa è stato fatto},
\textbf{dove si trova ogni componente}, \textbf{come funziona il flusso di
controllo}, \textbf{come si aggiungono nuove strategie} e \textbf{quali sono
gli step futuri e le scelte aperte}. L'obiettivo è metterti in condizione di
leggere il codice con una mappa in mano e di decidere quali migliorie
sperimentare.

Convenzione: tutti i percorsi sono relativi alla radice del progetto
\file{asa-deliveroo-agent/}. I termini del corso (belief, desire, intention,
plan) sono lasciati in inglese perché corrispondono ai nomi nel codice.

% ===========================================================================
\section{Il progetto in sintesi}
% ===========================================================================

\subsection{Requisiti d'esame coperti}

Il progetto realizza l'infrastruttura per i quattro pilastri dell'esame
(pesi di valutazione tra parentesi):

\begin{itemize}
  \item \textbf{Agente A — BDI (30\%)}: agente con revisione di belief e
    intention, strategia di gioco selezionabile, validato su scenari
    Challenge~1. Entry point: \file{src/main-bdi.js}.
  \item \textbf{Agente B — LLM + coordinamento (30\%)}: agente che
    interpreta le richieste degli agenti-missione, adatta la strategia e
    coordina con l'Agente A. Entry point: \file{src/main-llm.js}.
  \item \textbf{PDDL (20\%)}: planner PDDL integrato come membro della plan
    library per l'intention \code{go\_to}, opzionale a runtime.
    File: \file{src/planning/PddlPlanner.js}.
  \item \textbf{Report + presentazione (20\%)}: scheletro LaTeX in
    \file{report/} con commenti-guida per ogni sezione.
\end{itemize}

\subsection{Stato attuale (validato)}

Il software è \textbf{funzionante e testato}:

\begin{itemize}
  \item \textbf{Smoke test offline} (\code{npm test}): assertion su grafo,
    pathfinding, revisione belief, strategie, parsing missioni, generazione
    problemi PDDL e normalizzazione ack. Non richiede server né rete.
  \item \textbf{Test live} (11 giugno 2026, server locale, mappa
    \code{empty\_10}, 60 secondi): la strategia \code{reward-distance} ha
    totalizzato \textbf{800 punti con 30 parcel consegnati}. Il primo test
    (punteggio 173) ha rivelato un bug reale poi corretto
    (sezione~\ref{sec:bug-ack}).
\end{itemize}

% ===========================================================================
\section{Architettura generale}
% ===========================================================================

\subsection{Il ciclo di controllo BDI}

L'architettura segue il modello dei laboratori del corso: \emph{osserva
$\rightarrow$ rivedi i belief $\rightarrow$ genera opzioni $\rightarrow$
seleziona/rivedi l'intention $\rightarrow$ pianifica $\rightarrow$ agisci}.

\begin{verbatim}
        eventi socket (map, config, you, sensing, msg)
                          |
                   +------v------+
                   | BeliefBase  |  grafo mappa, io, parcel, agenti,
                   +------+------+  config, missione, compagno, claim
                          |
        OptionGenerator   (che cosa È POSSIBILE fare)
                          |
        Strategy          (che cosa CONVIENE fare - intercambiabile)
                          |                      ^
        IntentionRevision (commit/replace        |  l'LLM scrive qui:
                           con isteresi)         |  vincoli e goal di
                          |                      |  missione nei belief
        PlanLibrary       (COME farlo: GoPickUp, DeliverCarried,
                           GoToMissionTarget, Explore, Wait,
                           PddlGoTo -> FollowPathGoTo)
                          |
        ActionExecutor    (move/pickup/putdown serializzati,
                           gate "semaforo rosso")
\end{verbatim}

\subsection{Principi di design (da rispettare nelle evoluzioni)}

\begin{enumerate}
  \item \textbf{Infrastruttura vs strategia}: le strategie \emph{decidono}
    (ritornano l'opzione migliore), l'infrastruttura \emph{esegue}
    (pianifica, muove, valida, logga). Mai mescolare i due livelli.
  \item \textbf{Movimento deterministico}: pathfinding BFS sul grafo,
    nessuna decisione LLM nel ciclo passo-passo.
  \item \textbf{LLM ad alto livello}: interpreta le missioni e produce
    \emph{strutture dati validate} che entrano nei belief; non chiama mai
    direttamente le API di gioco.
  \item \textbf{Azioni serializzate}: il server rifiuta (con penalità)
    un'azione inviata mentre la precedente è in corso; ogni azione passa
    da un'unica coda di promise.
  \item \textbf{Niente costanti magiche}: tutti i parametri stanno in
    \file{src/config.js} o nelle opzioni dei costruttori delle strategie.
  \item \textbf{Mai fidarsi di una sola forma di ack}: i server del corso
    differiscono dai tipi dichiarati nell'SDK (verificato dal vivo).
\end{enumerate}

% ===========================================================================
\section{Mappa del repository}
% ===========================================================================

\begin{longtable}{@{}p{0.42\textwidth}p{0.52\textwidth}@{}}
\toprule
\textbf{Percorso} & \textbf{Responsabilità} \\
\midrule
\file{src/main-bdi.js} & Entry point Agente A; costruisce e avvia l'intero runtime (funzione \code{buildAgentRuntime} riusata anche dall'Agente B). \\
\file{src/main-llm.js} & Entry point Agente B; runtime BDI + livello LLM (interprete missioni, risposta alle domande, inoltro al compagno). \\
\file{src/config.js} & Configurazione runtime centralizzata (env + default documentati). \\
\midrule
\file{src/core/BeliefBase.js} & Stato del mondo e revisione dei belief; stato missione; riconciliazione. \\
\file{src/core/OptionGenerator.js} & Genera le opzioni candidate e i criteri di validità. \\
\file{src/core/Intention.js} & Un'opzione a cui l'agente si è impegnato; prova i piani in ordine. \\
\file{src/core/IntentionRevision.js} & Politica replace con isteresi; loop che consuma le intention. \\
\file{src/core/PlanLibrary.js} & Tutti i piani eseguibili + assemblaggio della libreria di default. \\
\file{src/core/ActionExecutor.js} & Serializzazione azioni; gate di movimento (red light). \\
\file{src/core/AgentLoop.js} & Collante: eventi $\rightarrow$ belief $\rightarrow$ deliberazione. \\
\midrule
\file{src/strategies/StrategyBase.js} & Interfaccia strategia (\code{utility}, \code{selectOption}). \\
\file{src/strategies/*Strategy.js} & Le 4 strategie esistenti (sez.~\ref{sec:strategie}). \\
\file{src/strategies/index.js} & Registro: id $\rightarrow$ classe; \code{createStrategy()}. \\
\midrule
\file{src/planning/GridGraph.js} & Grafo orientato della mappa; frecce; blocchi statici e temporanei. \\
\file{src/planning/PathPlanner.js} & BFS; distanze da me; distanza-consegna precalcolata; helper di scoring. \\
\file{src/planning/PddlPlanner.js} & Generazione problema PDDL dai belief + solver online; logging. \\
\file{src/planning/pddlDomain.js} & Dominio STRIPS leggibile + mappa azione$\rightarrow$direzione. \\
\midrule
\file{src/communication/MessageTypes.js} & Busta del protocollo \code{asa-team-v1} e tipi di messaggio. \\
\file{src/communication/TeamProtocol.js} & Handshake, heartbeat, claim, mission-update, ack. \\
\midrule
\file{src/llm/LlmClient.js} & Wrapper endpoint OpenAI-compatibile (LiteLLM del corso); lazy. \\
\file{src/llm/MissionInterpreter.js} & Messaggio missione $\rightarrow$ oggetto strutturato (LLM o fallback regex). \\
\file{src/llm/prompts.js} & Prompt di sistema (interpretazione missioni; tool-agent futuro). \\
\file{src/llm/tools.js} & Registro tool validati ad alto livello (pattern lab8). \\
\midrule
\file{src/metrics/MetricsCollector.js} & Contatori e timeline punteggio di una run. \\
\file{src/metrics/RunLogger.js} & Log JSON-lines per run + scrittura riassunto risultati. \\
\midrule
\file{src/utils/} & \code{manhattan}, \code{sleep}, chiavi tile, conversione clock-event, \code{normalizeIdList}, estrazione JSON. \\
\midrule
\file{scripts/run-bdi.js, run-llm.js} & Lancio agenti con flag CLI che sovrascrivono l'env. \\
\file{scripts/run-experiment.js} & Run a tempo: avvia, gioca, scrive risultati, esce. \\
\file{scripts/smoke-test.js} & Test offline completo (\code{npm test}). \\
\midrule
\file{experiments/} & \file{logs/} (JSON-lines) e \file{results/} (riassunti); entrambi git-ignored. \\
\file{report/} & Report d'esame (scheletro) + questo documento. \\
\file{notebooks/} & Solo analisi offline dei risultati (mai runtime). \\
\file{CLAUDE.md} & Regole del progetto, inclusa la regola ``la documentazione segue il codice''. \\
\bottomrule
\end{longtable}

% ===========================================================================
\section{Il nucleo BDI in dettaglio}
% ===========================================================================

\subsection{BeliefBase — lo stato del mondo}

\file{src/core/BeliefBase.js} è l'unica fonte di verità. Campi principali:
\code{me} (posizione/punteggio, autoritativi solo via evento \code{you} e
ack dei move), \code{graph} (sez.~\ref{sec:grafo}), \code{parcels},
\code{agents}, \code{tileLastSeen}, \code{config} (config di gioco del
server: mai hardcodare!), \code{mission}, \code{teammate}, \code{claims}.

Principi di revisione implementati (dal lab2 e dalla knowledge base):

\begin{itemize}
  \item \textbf{Timestamp su tutto}: ogni belief dinamico ha
    \code{lastSeen}.
  \item \textbf{Proiezione del decay}: la reward ricordata è un limite
    superiore; \code{projectedReward()} sottrae un punto per ogni
    \emph{tick intero} di decay trascorso da \code{lastSeen} (semantica a
    scatti del server; fix di precisione dell'11/06).
  \item \textbf{Evidenza negativa}: un parcel ricordato viene cancellato
    solo quando il suo tile \emph{rientra nel campo visivo} e il parcel
    non c'è più. L'assenza di sensing non è evidenza di assenza.
  \item \textbf{Riconciliazione} (``la mia conoscenza era sbagliata, non il
    mondo''): \code{clearCarried()} dimentica i parcel che credevo di
    trasportare quando il server mi smentisce; \code{markTilePickedUp()}
    marca come trasportati i parcel del mio tile quando l'ack di pickup
    non contiene id utilizzabili.
\end{itemize}

Lo stato missione (\code{mission}) è il canale con cui l'LLM influenza il
BDI: \code{setMission()} traduce una missione strutturata in effetti
concreti — tile vietati bloccati nel grafo, tile di consegna esclusi,
politiche di putdown (\code{deliverExactly}, \code{deliverMaxValue}), gate
di movimento (\code{movementAllowed}), goal posizionali
(\code{mission.active}).

\subsection{OptionGenerator — che cosa è possibile}

Genera le opzioni candidate ad ogni deliberazione:
\code{go\_pick\_up} (un parcel libero, non reclamato dal compagno, con
reward proiettata $>0$), \code{deliver\_carried} (se trasporto qualcosa),
\code{go\_to\_mission\_target} (se c'è una missione attiva con coordinate),
più i fallback \code{explore} e \code{wait}. Espone anche
\code{isStillValid(option, beliefs)}: le condizioni di abbandono
(\emph{achieved / impossible / not worthwhile}) usate dal loop di revisione.

\textbf{Dove intervenire}: se una nuova strategia ha bisogno di un nuovo
tipo di obiettivo (es.\ ``camping su uno spawner''), il tipo di opzione si
aggiunge qui, il piano corrispondente in \file{PlanLibrary.js}, e la
strategia gli assegna un'utilità.

\subsection{IntentionRevision — impegno con isteresi}

Politica \emph{replace} raffinata: l'agente è impegnato su \textbf{una}
intention alla volta; una nuova opzione la sostituisce solo se
\[
  \mathit{utility}_{\text{nuova}} > \mathit{utility}_{\text{corrente}} + \mathit{margine}
\]
con margine additivo configurabile (\code{agent.hysteresisMargin}, default
5 $\approx$ 5 punti di reward). Questo evita il \emph{thrashing} fra due
parcel di valore quasi identico. Ripresentare la stessa opzione aggiorna
solo l'utilità (per confronti freschi) senza riavviare l'intention.
\code{abortCurrent()} interrompe un'intention diventata invalida (il check
lo fa \code{AgentLoop.deliberate()} ad ogni ciclo).

\subsection{Intention e PlanLibrary — dal volere al fare}

Una \code{Intention} prova \emph{in ordine} tutte le classi di piano
applicabili; se un piano fallisce (lancia un'eccezione) si passa al
successivo; se falliscono tutti, l'intention fallisce e la deliberazione
successiva riparte da capo. È il meccanismo con cui ``i fallimenti vengono
riportati al livello delle intention''. Piani esistenti:

\begin{itemize}
  \item \code{GoPickUp}: sub-intention \code{go\_to} + \code{pickup};
    ack vuoto $\Rightarrow$ gara persa, belief cancellato, fallimento.
  \item \code{DeliverCarried}: consegna al delivery più vicino; selezione
    \emph{mission-aware} dei parcel da posare (esattamente $N$ / valore
    $\le$ soglia); riconciliazione su ack vuoto.
  \item \code{GoToMissionTarget}: raggiunge il target di missione più
    vicino; \code{deliver\_at} posa i parcel; \code{holdAtTarget} attende
    (placeholder 5\,s, da sostituire con sincronizzazione esplicita).
  \item \code{Explore}: va allo spawner visto meno di recente (massimizza
    informazione per tile percorso).
  \item \code{Wait}: pausa breve (ultima spiaggia).
  \item \code{PddlGoTo} $\rightarrow$ \code{FollowPathGoTo}: registrati in
    quest'ordine; se il PDDL è disabilitato o fallisce, il BFS subentra
    automaticamente.
\end{itemize}

\code{FollowPathGoTo} è il cuore deterministico: ricalcola il path BFS,
esegue un passo alla volta, e su un move fallito fa \code{moveRetries}
tentativi (l'ostacolo dinamico può spostarsi); se il blocco persiste,
marca il tile come \emph{soft-blocked} per \code{softBlockMs} millisecondi
(il prossimo BFS ci gira attorno) e fallisce il piano.

\subsection{ActionExecutor — serializzazione e semaforo}

Tutte le azioni di gioco passano da un'unica catena di promise: la
successiva parte solo dopo l'ack della precedente (il server rifiuta con
penalità le azioni concorrenti — \emph{ActionMutex}). Il
\code{movementGate} è un predicato sostituibile: quando la missione
``red light'' chiude il gate, \code{move()} \emph{attende} invece di
fallire — l'intention si congela e riprende al verde, che è esattamente il
comportamento premiato.

\subsection{AgentLoop — il collante}

Collega gli eventi socket ai belief, attende mappa e identità, avvia il
loop di revisione e delibera: ad ogni evento \code{you}/\code{sensing},
dopo ogni intention conclusa, e comunque ogni
\code{deliberationIntervalMs} (250\,ms) — necessario perché il sensing
arriva \emph{solo quando qualcosa cambia}: il silenzio non deve congelare
l'agente.

% ===========================================================================
\section{Navigazione deterministica}\label{sec:grafo}
% ===========================================================================

\subsection{GridGraph — la mappa come grafo orientato}

Costruito una volta dall'evento \code{map}. Punti chiave:

\begin{itemize}
  \item \textbf{Digrafo dal primo giorno}: i tile freccia
    (\code{↑→↓←}) vietano l'ingresso \emph{contro} la freccia; la regola è
    codificata omettendo l'arco. Le mappe Challenge~1 n.~4--6 (circuiti a
    senso unico) funzionano senza modifiche.
  \item \textbf{Blocchi statici}: \code{blockTiles()} per i vincoli di
    missione (``non passare per (x,y)''); \code{setForbiddenDeliveries()}
    esclude tile di consegna vietati.
  \item \textbf{Soft block}: tile temporaneamente evitati (agente fermo
    sul percorso) con scadenza automatica.
  \item \textbf{Distanza-consegna precalcolata}: BFS multi-sorgente dai
    delivery sugli \emph{archi invertiti} (corretto sul digrafo: con i
    sensi unici la distanza \emph{verso} un delivery non è simmetrica).
    Lookup $O(1)$ usato dalle utilità.
\end{itemize}

\subsection{PathPlanner — BFS e helper di scoring}

\code{shortestPath} (BFS, costi unitari), \code{distancesFrom} (una BFS
sola per prezzare \emph{tutte} le opzioni di una deliberazione),
\code{nearestDelivery}, e \code{scoringHelpers()} che fornisce alle
strategie: \code{distanceTo(x,y)}, \code{deliveryDistanceFrom(x,y)},
\code{decayPerTile} (reward persa per passo per parcel trasportato
$= \mathit{movement\_duration} / \mathit{decay\_interval}$).

% ===========================================================================
\section{PDDL}
% ===========================================================================

Scelta architetturale (motivata nella knowledge base): il PDDL è un
\textbf{membro della plan library}, non il loop esterno. Il mondo cambia ad
ogni frame e il solver online ha latenze di secondi: ripianificare via HTTP
ad ogni tick non è praticabile; servire l'intention \code{go\_to} con un
piano STRIPS sì, ed è un'integrazione ``significativa'' richiesta
dall'esame.

\begin{itemize}
  \item \file{pddlDomain.js}: dominio leggibile con predicati
    \code{(at ?t)} e archi direzionali \code{(up/down/left/right ?from ?to)};
    4 azioni \code{move-*}. I sensi unici sono codificati gratis:
    l'ingresso vietato semplicemente non ha il predicato-arco.
  \item \file{PddlPlanner.js}: genera il problema dai belief (regione
    raggiungibile via BFS, con limite \code{pddl.maxTiles}), chiama
    \code{onlineSolver} di \code{@unitn-asa/pddl-client} (caricato lazy:
    senza pacchetto o con \code{PDDL\_ENABLED=false} non se ne paga il
    costo), converte i passi in direzioni, \textbf{logga ogni chiamata e
    ogni fallimento} (\code{pddl\_plan} / \code{pddl\_failure}) — dati
    pronti per il confronto BFS-vs-PDDL nel report.
\end{itemize}

\textbf{Evoluzione prevista}: estendere il dominio con azioni
\code{pickup}/\code{putdown} (pianificare intere sequenze
raccogli-e-consegna) o con la spinta delle casse
(\code{domain-deliveroojs-crates.pddl} come base) per le mappe di pratica
sokoban-like.

% ===========================================================================
\section{Le strategie}\label{sec:strategie}
% ===========================================================================

\subsection{L'interfaccia}

\file{StrategyBase.js}: una strategia implementa
\code{utility(option, beliefs, helpers)} (e/o sovrascrive
\code{selectOption} per logiche non a utilità). La selezione di default è
l'argmax delle utilità; le utilità vengono scritte \emph{sulle opzioni}
perché la revisione le usa per l'isteresi. Convenzione di scala: ``punti di
reward attesi'', così opzioni di tipo diverso restano confrontabili. La
\code{StrategyBase} dà già a \emph{tutte} le strategie il rispetto delle
missioni posizionali (i bonus $\pm200\ldots\pm1000$ dominano il reddito da
parcel).

\subsection{Le quattro strategie esistenti}

\begin{longtable}{@{}p{0.25\textwidth}p{0.69\textwidth}@{}}
\toprule
\textbf{Id} & \textbf{Idea e formula} \\
\midrule
\code{greedy-nearest} & Baseline di controllo: insegue il parcel
  raggiungibile più vicino ($u = 1000 - d$); consegna solo quando non c'è
  altro da raccogliere ($u = 500 - d_{\text{delivery}}$). Ignora valore e
  decay di proposito. \\
\midrule
\code{reward-distance} & Massimizza il valore \emph{consegnato}:
  $u_{\text{pick}} = \mathit{proj}(p) - (d_{me\to p} + d_{p\to del})\cdot
  \mathit{decay} \cdot (c+1)$;
  $u_{\text{del}} = \sum \mathit{proj}(carried) - d_{me\to del}\cdot
  \mathit{decay}\cdot c - \mathit{deliverBias}$, dove $c$ = parcel
  trasportati e \code{deliverBias} (default 10) incoraggia il batching. \\
\midrule
\code{delivery-threshold} & Come sopra, ma accumula fino a soglia
  (\code{minCarried}, default 3, o \code{minValue}) e poi la consegna
  domina (+1000). Pensata per mappe dove le consegne sono il collo di
  bottiglia (es.\ 26c1\_5: 91 spawner, 4 delivery). \\
\midrule
\code{mission-aware} & \code{reward-distance} + obbedienza allo stato
  missione: peso configurabile sui goal di missione
  (\code{missionWeight}); boost (+200) alla consegna quando la politica
  \code{deliverExactly}/\code{deliverMaxValue} è soddisfacibile. Default
  dell'Agente B. \\
\bottomrule
\end{longtable}

\subsection{Come aggiungere una nuova strategia (passo-passo)}

\begin{enumerate}
  \item Crea \file{src/strategies/MiaStrategia.js}:
\begin{lstlisting}
import { StrategyBase } from './StrategyBase.js';
// oppure: extends RewardDistanceStrategy per ereditarne le utilita'

export class MiaStrategia extends StrategyBase {
  static id = 'mia-strategia';

  constructor(options = {}) {
    super(options);
    this.mioParametro = options.mioParametro ?? 42; // niente costanti magiche
  }

  utility(option, beliefs, helpers) {
    if (option.type === 'go_pick_up') {
      const parcel = beliefs.parcels.get(option.parcelId);
      const d = helpers.distanceTo(option.x, option.y);   // distanza VERA (BFS)
      if (!Number.isFinite(d)) return -Infinity;          // irraggiungibile
      return beliefs.projectedReward(parcel) - d * helpers.decayPerTile;
    }
    return super.utility(option, beliefs, helpers); // missioni/explore/wait
  }
}
\end{lstlisting}
  \item Registrala in \file{src/strategies/index.js}: un import + una riga
    nella lista \code{STRATEGY\_CLASSES}.
  \item Lanciala: \code{node scripts/run-bdi.js --strategy mia-strategia}
    (o \code{STRATEGY=mia-strategia} nel \file{.env}).
  \item Confrontala: \code{node scripts/run-experiment.js --strategy
    mia-strategia --duration 180 --label 26c1\_3}, almeno 5 run per
    coppia (mappa, strategia), confronto sulle medie.
\end{enumerate}

Strumenti a disposizione dentro \code{utility}: tutti i belief
(\code{beliefs.parcels}, \code{beliefs.agents}, \code{beliefs.carried()},
\code{beliefs.projectedReward(p)}, \code{beliefs.mission},
\code{beliefs.teammate}, \code{beliefs.claims}) e gli helper geometrici
(\code{distanceTo}, \code{deliveryDistanceFrom}, \code{decayPerTile}).
Per logiche non a punteggio (es.\ macchine a stati) sovrascrivi
\code{selectOption} e ritorna direttamente l'opzione scelta (ricordati di
assegnarle \code{option.utility}, serve all'isteresi).

% ===========================================================================
\section{L'Agente B: livello LLM}
% ===========================================================================

\subsection{Flusso}

\file{main-llm.js} costruisce \emph{lo stesso} runtime BDI dell'Agente A
(deve pur sempre muoversi e consegnare) e ci aggiunge un listener sui
messaggi non-protocollo: ogni messaggio (tipicamente lo \code{shout} di un
agente-missione) passa da \code{MissionInterpreter.interpret()}; il
risultato strutturato viene applicato ai propri belief
(\code{setMission}) e inoltrato all'Agente A via
\code{protocol.sendMissionUpdate()}. Le missioni \code{question\_answer}
(es.\ ``Calcola $5*(5+3)/2$'') vengono risposte direttamente in chat senza
muoversi (livello~1 della Challenge~2).

\subsection{MissionInterpreter — doppio binario}

\begin{itemize}
  \item \textbf{Binario LLM} (se configurato in \file{.env}): un prompt
    (\file{prompts.js}) estrae kind, coordinate, segno del bonus, soglie;
    l'output JSON è \textbf{validato} (\code{\#validate}) prima di toccare
    i belief; output invalido $\Rightarrow$ fallback.
  \item \textbf{Binario deterministico} (sempre disponibile): regex sul
    catalogo missioni della Challenge~2 — \code{go\_to},
    \code{deliver\_at} (anche in forma vietata), \code{question\_answer}
    (con valutatore aritmetico sanificato), \code{deliver\_exactly\_n},
    \code{deliver\_less\_value\_than}, \code{one\_pickup\_another\_deliver},
    \code{red\_light\_green\_light}.
  \item \textbf{Eccezione di latenza}: i cambi di stato
    \code{RED LIGHT}/\code{GREEN LIGHT} sono \emph{sempre} parsati senza
    LLM (un round-trip di rete durante il rosso costerebbe penalità).
\end{itemize}

\subsection{Tools (per esperimenti futuri)}

\file{tools.js} contiene il registro in stile lab8 (\code{get\_state},
\code{set\_mission}, \code{block\_tiles}, \code{set\_movement\_allowed},
\code{send\_to\_teammate}, \code{reply\_to\_agent}) con dispatch validato
(\code{executeTool}). Non è nel flusso di default: serve a sperimentare un
loop in cui l'LLM orchestri i tool per richieste aperte. Volutamente
\textbf{non esiste} un tool ``muoviti di un passo''.

% ===========================================================================
\section{Coordinamento A--B}
% ===========================================================================

Protocollo \code{asa-team-v1} (\file{MessageTypes.js}): busta JSON
\code{\{protocol, type, payload, from, ts\}}. I messaggi non-busta (es.\
shout delle missioni) vengono ignorati dal protocollo e gestiti dal
livello LLM.

\begin{longtable}{@{}llp{0.55\textwidth}@{}}
\toprule
\textbf{Tipo} & \textbf{Payload} & \textbf{Uso} \\
\midrule
\code{hello} & \code{\{name, isReply?\}} & Scoperta del compagno: shout
  all'avvio; chi riconosce il nome atteso (\code{TEAMMATE\_NAME}) registra
  l'id e risponde via \code{say}. \\
\code{position} & \code{\{x, y, carrying, intention\}} & Heartbeat
  $\sim$1\,s (raddoppia l'osservazione effettiva del team). \\
\code{intention} & \code{\{key, type\}} & Annuncio intention corrente. \\
\code{claim} & \code{\{parcelId\}} & Deconflitto target: primo che reclama
  vince; l'\code{OptionGenerator} dell'altro scarta il parcel. \\
\code{mission-update} & \code{\{mission\}} & B $\rightarrow$ A: missione
  strutturata da applicare ai belief. \\
\code{request-help} & \code{\{reason, x?, y?\}} & Richiesta d'aiuto
  (gestione strategica: TODO). \\
\code{handover} & \code{\{parcelId, x, y, role\}} & Passaggio di parcel
  (26c2\_8) — coreografia completa: TODO. \\
\code{ack} & \code{\{about\}} & Conferma; via callback di \code{ask} se
  presente. \\
\bottomrule
\end{longtable}

Nota operativa: \code{say} per i flussi normali; \code{ask} (bloccante,
timeout server 1\,s) riservato ai punti di sincronizzazione, con timeout
proprio lato client.

% ===========================================================================
\section{Metriche, log ed esperimenti}
% ===========================================================================

\begin{itemize}
  \item \textbf{RunLogger}: un file JSON-lines per agente per run in
    \file{experiments/logs/} (\code{\{t, event, ...\}}). Eventi:
    punteggio, pickup/consegne (con \code{count} e \code{ids}
    normalizzati), intention avviate/concluse/fallite/abortite, piani
    falliti, chiamate/fallimenti PDDL, interpretazioni LLM, messaggi.
  \item \textbf{MetricsCollector}: contatori in memoria + timeline del
    punteggio; \code{summary()} scritto in \file{experiments/results/} a
    fine run.
  \item \textbf{run-experiment.js}: run a tempo riproducibile:
    \code{--agent bdi|llm --strategy <id> --duration <s> --label <scenario>}.
  \item \textbf{Metodologia}: gli spawn sono casuali $\Rightarrow$ almeno
    5 run per coppia (scenario, strategia) e confronto su media$\pm$dev.\
    standard (stesso pattern del benchmark del professore).
  \item Gli output grezzi sono \textbf{git-ignored}: i numeri che servono
    al report vanno copiati nei sorgenti del report o nei README.
\end{itemize}

% ===========================================================================
\section{Storia della validazione (e lezioni)}\label{sec:bug-ack}
% ===========================================================================

\subsection{Il bug dei ``parcel fantasma''}

Primo test live (60\,s, \code{empty\_10}): punteggio 173, ma 352 cambi di
intention e 336 putdown falliti. Diagnosi dai log: gli \textbf{ack del
server non contenevano il campo \code{id}} (forma diversa dai tipi
dichiarati nell'SDK). Risultato: dopo ogni consegna il belief ``sto
trasportando X'' restava vivo (parcel fantasma) e l'agente riprovava il
putdown all'infinito finché la proiezione del decay non azzerava il
fantasma.

\textbf{Fix} (tre pezzi, tutti in chiave ``belief reconciliation''):
\code{normalizeIdList()} accetta ack come oggetti \code{\{id\}}, stringhe
o forme miste; \code{markTilePickedUp()} come fallback del pickup;
\code{clearCarried()} quando il server smentisce il carry. Secondo test:
\textbf{punteggio 800, 30 consegne, zero retry futili}. La coppia di run è
un esempio già pronto, con dati, del valore della revisione dei belief per
il report (sezione belief revision).

\subsection{Fix del decay (contributo Codex)}

\code{projectedReward} sottraeva frazioni di tick
($\lfloor r - t/T \rfloor$); il server decrementa a scatti interi. Ora:
$r - \lfloor t/T \rfloor$, con guardia \code{Math.max(0, elapsed)} contro
skew d'orologio. Pinnato da assertion nello smoke test.

\subsection{Lezione generale}

\emph{Mai fidarsi dei tipi dichiarati}: i server del corso (locale vs
challenge) possono differire dall'SDK. Ogni dato che entra nei belief passa
da normalizzazione e, dove possibile, da riconciliazione con la realtà
osservata.

% ===========================================================================
\section{Step futuri}
% ===========================================================================

In ordine di priorità suggerito:

\begin{enumerate}
  \item \textbf{Tuning e validazione Challenge 1} (pilastro da 30\%):
    girare le 8 mappe \code{26c1\_*}; tarare \code{deliverBias},
    \code{minCarried}, \code{hysteresisMargin} per mappa; verificare i
    circuiti a senso unico (4--5) e la competizione con l'NPC onnisciente.
  \item \textbf{Utilità opponent-aware}: se un avversario visibile è più
    vicino del nostro agente al parcel target, abbandonarlo (l'NPC
    ``intelligent'' è onnisciente: la gara si vince solo in velocità).
  \item \textbf{Coreografia handover (26c2\_8)}: negoziare il tile di
    rendezvous via \code{ask}, sequenziare putdown-prima-del-pickup
    (due agenti non condividono il tile), scambiarsi gli id dei parcel.
  \item \textbf{Go-to-and-wait di squadra (26c2\_10)}: sostituire
    l'attesa fissa di 5\,s con sincronizzazione esplicita
    (\code{position} + \code{ack}).
  \item \textbf{PDDL esteso}: azioni pickup/putdown nel dominio;
    esperimento BFS-vs-PDDL (latenza, lunghezza piano, fallimenti) per la
    sezione PDDL del report.
  \item \textbf{Tool loop LLM} (opzionale): l'LLM orchestra i tool di
    \file{tools.js} per richieste aperte (pattern lab8 07).
  \item \textbf{Esperimenti sistematici + report}: $\ge 5$ run per
    (mappa, strategia); notebook di analisi; riempire le sezioni del
    report seguendo i commenti-guida.
\end{enumerate}

% ===========================================================================
\section{Scelte aperte e parametri da esplorare}
% ===========================================================================

\subsection{Manopole esistenti (nessuna modifica al codice richiesta)}

\begin{longtable}{@{}llp{0.52\textwidth}@{}}
\toprule
\textbf{Parametro} & \textbf{Default} & \textbf{Che cosa governa / che cosa provare} \\
\midrule
\code{hysteresisMargin} & 5 & Stabilità dell'impegno: alzarlo riduce i
  cambi di target, abbassarlo rende l'agente più opportunista. Confrontare
  i \code{intentionChanges} nei risultati. \\
\code{deliverBias} & 10 & Propensione al batching di
  \code{reward-distance}: su mappe con decay veloce conviene abbassarlo. \\
\code{minCarried} / \code{minValue} & 3 / $\infty$ & Soglie di
  \code{delivery-threshold}: idealmente derivarle dal config dello
  scenario (\code{parcels.max}, velocità di decay) invece che fissarle. \\
\code{missionWeight} & 1 & Aggressività sulle missioni di
  \code{mission-aware}. \\
\code{deliberationIntervalMs} & 250 & Frequenza di riconsiderazione:
  scenari dinamici $\rightarrow$ più alta; stabili $\rightarrow$ più
  bassa (meno overhead). \\
\code{moveRetries} / \code{softBlockMs} & 2 / 3000 & Pazienza contro
  ostacoli dinamici vs velocità di re-routing. \\
\code{pddl.enabled} / \code{maxTiles} & false / 1600 & Attivazione PDDL e
  limite dimensione problema. \\
\bottomrule
\end{longtable}

\subsection{Idee strategiche da prototipare}

\begin{itemize}
  \item \textbf{Partizionamento della mappa} per il team (regioni di
    Voronoi sui delivery) per non farsi concorrenza — da misurare contro
    il semplice claim.
  \item \textbf{Camping sugli spawner} in mappe a spawn lento: piazzarsi
    al centroide degli spawner invece di esplorare.
  \item \textbf{Stima del movimento avversario} dalle posizioni
    successive nel sensing (i belief \code{agents} hanno già
    \code{lastSeen}).
  \item \textbf{Soglie auto-adattive}: leggere \code{parcels.max},
    \code{generation\_event}, \code{decaying\_event} dal config e derivare
    i parametri di batching per scenario.
  \item \textbf{Sfruttare il throttling degli spawn}: i parcel trasportati
    contano nel limite \code{parcels.max} $\Rightarrow$ accumulare affama
    la mappa (rilevante in 26c1\_4/5 con max=5).
\end{itemize}

\subsection{Domande aperte da verificare a runtime}

Dalla knowledge base (\file{context/game\_knowledge/06}): la capacità è
davvero non applicata ai player sul server delle challenge? Il sensing
include i parcel trasportati da me? (localmente: sì, ma via ack ci
arriviamo comunque). Latenza fra evento-missione e accredito del bonus?
Versione del server delle challenge vs locale? Conviene un probe-agent che
logga tutti gli eventi grezzi per 5 minuti per mappa.

% ===========================================================================
\section{Regole di lavoro e comandi rapidi}
% ===========================================================================

\subsection{La regola ``la documentazione segue il codice''}

Codificata in \file{CLAUDE.md} e \textbf{obbligatoria}: ogni modifica al
comportamento del codice aggiorna nello stesso cambiamento (1) il
\file{README.md} (implementato/assunzioni), (2) \file{scripts/smoke-test.js}
(assertion sul nuovo comportamento; \code{npm test} deve passare offline),
(3) \file{experiments/README.md} se cambiano eventi/contatori/formati,
(4) \file{notebooks/README.md} se cambia il formato dei risultati,
(5) i commenti-guida in \file{report/sections/*.tex}, (6) i commenti JSDoc.
Una modifica non è completa finché la checklist non è stata percorsa.

\subsection{Comandi}

\begin{lstlisting}[language={}]
npm test                                        # smoke test offline
node scripts/run-bdi.js --strategy <id>         # Agente A
node scripts/run-llm.js  --name agentB          # Agente B
node scripts/run-experiment.js --strategy <id> \
     --duration 180 --label <scenario>          # run a tempo con risultati
\end{lstlisting}

Variabili d'ambiente principali (vedi \file{.env.example} commentato):
\code{HOST}, \code{TOKEN} (salvare quello stampato al primo avvio per
mantenere l'identità), \code{NAME}, \code{TEAMMATE\_NAME},
\code{STRATEGY}, \code{PDDL\_ENABLED}, \code{LITELLM\_*}.

\bigskip
\noindent\textit{Fine del documento. Aggiornarlo quando l'architettura
evolve: fa parte della documentazione coperta dalla regola di cui sopra.}

\end{document}
